\section{问题描述与控制方程}

\subsection{一维 Riemann 问题}

\subsubsection{计算目标}

本项目采用 MacCormack 格式求解一维 Riemann 问题的数值解，并将数值解与
半解析解（精确解）进行对比。具体研究目标如下。

\begin{enumerate}[leftmargin=2em]
    \item 用 MacCormack（带人工粘性版） 求解7个1D-Riemann初值问题：
    \begin{itemize}
        \item Sod
        \item LAX
        \item 亚声速后退
        \item Sjogreen supersonic Expansion超声速后退
        \item 接触间断 - 双膨胀激波
        \item 双激波间断
        \item 纯接触间断问题 - 注意到反而是有些高级格式在这个问题上比不过MacCormark格式，可以拓展对比一下）=> 不同格式优秀的点还是有差异的
    \end{itemize}
    \item 将 MacCormack 数值解与精确解进行对比；
    \item 讨论人工粘性的有无对解的影响；
    \item 研究人工粘性经验参数 $\beta$ 对数值解的影响；
    \item 研究 CFL 数对数值解的影响，并讨论人工粘性系数与 CFL 数之间
    并非简单单调关系的现象；
    \item 先以密度作为人工粘性开关函数，并进一步比较采用压力或速度作为
    开关函数时的差异；
    \item 在不加入人工粘性的情况下，通过加密网格考察色散效应是否得到改善，
    并分析 MacCormack 格式固有色散误差与网格分辨率之间的关系。
\end{enumerate}

\subsubsection{Project 2 与 Project 0 精确解之间的关系}
Project 0 给出同一 Riemann 问题的精确解程序。虽然程序内部在求解超越方程时仍要用数值迭代方法，
但它对应的是精确 Riemann 解（与DNS的内涵完全不一样），可以作为 Project 2 数值解的参考解。Project 2 采用 MacCormack 格式计算近似解，
再与 Project 0 的结果比较，用来检查误差、稳定性和间断位置。

最终输出和比较时主要使用原始量 \(\bm W=(\rho,u,p)^{\mathrm T}\)。
\begin{equation}
\bm{W}(x,0)=
\begin{cases}
(\rho_L,u_L,p_L)^{\mathrm T}, & x<x_0,\\
(\rho_R,u_R,p_R)^{\mathrm T}, & x>x_0.
\end{cases}
\end{equation}

\subsection{一维 Euler 方程}

\begin{equation}
\frac{\partial \bm{Q}}{\partial t}
+
\frac{\partial \bm{F}(\bm{Q})}{\partial x}
=0,
\end{equation}

\begin{equation}
\bm{Q}=
\begin{bmatrix}
\rho\\
\rho u\\
\rho E
\end{bmatrix},
\qquad
\bm{F}=
\begin{bmatrix}
\rho u\\
\rho u^2+p\\
u(\rho E+p)
\end{bmatrix}.
\end{equation}

总能量满足

\begin{equation}
E=\frac{p}{\rho(\gamma-1)}+\frac{u^2}{2}.
\end{equation}
